Traffic congestion in Vietnam's major cities is causing serious economic losses. According to figures from the Department of Transport of Ho Chi Minh City, the city loses approximately 6 billion USD  each year due to traffic congestion [1]. In Hanoi, congestion-related losses are estimated at 1 to 1.2 billion USD per year [2]. These figures account for wasted fuel costs, lost working hours, and the broader impact on economic productivity.
Expanding road infrastructure requires large capital investments and long construction periods, making it an impractical solution in the short term. A more realistic approach is to optimize how traffic is controlled within the existing network, particularly at signalized intersections — the points that directly regulate vehicle flow and have a major influence on the operational efficiency of the entire network.
In Vietnam, most intersections still operate under a fixed-time signal control scheme, in which the cycle length and phase durations are calculated using the Webster method and kept unchanged regardless of actual traffic flow [3], [4]. This approach is simple and stable, but it cannot adapt when traffic surges or incidents occur, causing queues to build up and spread to neighboring intersections.
Trong những năm gần đây, Học Tăng Cường (Reinforcement Learning — RL) đã nổi lên như một hướng tiếp cận đầy triển vọng cho bài toán này. Khác với các phương pháp dựa trên quy tắc cố định, một tác nhân RL học chiến lược điều khiển thông qua việc tương tác trực tiếp với môi trường giao thông và nhận phản hồi dưới dạng phần thưởng, từ đó tự khám phá ra chính sách điều khiển tối ưu. 
Trong nghiên cứu này, ba thuật toán Deep RL được lựa chọn để khảo sát: Deep Q-Network (DQN), Double DQN (DDQN) và Proximal Policy Optimization (PPO). Trong đó DQN và DDQN thuộc nhóm phương pháp dựa trên giá trị, còn PPO thuộc nhóm dựa trên chính sách

Các đóng góp chính của khóa luận bao gồm:

Xây dựng môi trường mô phỏng giao thông dựa trên SUMO và mô hình hóa bài toán điều khiển đèn tín hiệu dưới dạng Quá trình Quyết định Markov (MDP) 
Triển khai và so sánh ba thuật toán RL (PPO, DQN, DDQN) kết hợp với ba hàm phần thưởng khác nhau (dựa trên hàng chờ, áp lực, và thời gian chờ có giới hạn) trên cùng một bộ điều kiện thực nghiệm.
Đánh giá khả năng mở rộng của các thuật toán khi chuyển từ nút giao đơn lẻ sang mạng lưới 2×2 với kiến trúc Independent chia sẻ trọng số chính sách, đồng thời phân tích ảnh hưởng của hàm phần thưởng đến sự phối hợp giữa nhiều tác nhân.
Xác định cấu hình điều khiển hiệu quả nhất cho từng loại mạng giao thông, cung cấp cơ sở thực nghiệm cho việc lựa chọn thuật toán và hàm phần thưởng trong các ứng dụng điều khiển đèn tín hiệu thực tế.
